\documentclass[11pt]{article}
\usepackage[margin=1in]{geometry}
\usepackage{amsmath,amssymb}

\title{MATH 375: Probability Theory}
\author{Mariana Restrepo}
\date{Fall 2026}

\begin{document}
\maketitle

\begin{center}
\large Lecture 6
\end{center}

\tableofcontents

\section{Lecture 6: Independence and Repeated Trials}

\subsection{Independence and complements}

Events $E$ and $F$ are independent when
\[
P(E\cap F)=P(E)P(F).
\]

If $E$ and $F$ are independent, then $E$ and $F^c$ are also
independent. To see why, split $E$ into disjoint pieces:
\[
E=(E\cap F)\cup(E\cap F^c).
\]
Consequently,
\[
\begin{aligned}
P(E\cap F^c)
&=P(E)-P(E\cap F)\\
&=P(E)-P(E)P(F)\\
&=P(E)\bigl(1-P(F)\bigr)\\
&=P(E)P(F^c).
\end{aligned}
\]
This is exactly the condition for $E$ and $F^c$ to be independent.

\subsection{Independence of three events}

Three events $E$, $F$, and $G$ are \textbf{mutually independent}
when all four conditions hold:
\[
\begin{aligned}
P(E\cap F)&=P(E)P(F),\\
P(E\cap G)&=P(E)P(G),\\
P(F\cap G)&=P(F)P(G),\\
P(E\cap F\cap G)&=P(E)P(F)P(G).
\end{aligned}
\]
Checking only the three pairs is not enough to establish mutual
independence.

\subsection{Repeated independent trials}

Suppose each trial independently succeeds with probability $p$.
Then each trial fails with probability $1-p$.

\paragraph{At least one success in $n$ trials.}
The complement is ``all $n$ trials fail,'' with probability
$(1-p)^n$. Therefore,
\[
\boxed{
P(\text{at least one success})=1-(1-p)^n.
}
\]

\paragraph{Exactly $k$ successes in $n$ trials.}
First choose which $k$ positions contain successes. For any
particular choice of positions, independence gives probability
$p^k(1-p)^{n-k}$. Hence
\[
\boxed{
P(\text{exactly $k$ successes})
=\binom{n}{k}p^k(1-p)^{n-k}.
}
\]

\subsection{Which dice sum appears first?}

Independently roll two fair dice until either a sum of 5 or a
sum of 7 appears. What is the probability that 5 appears first?

In one roll,
\[
P(\text{sum 5})=\frac4{36},\qquad
P(\text{sum 7})=\frac6{36}.
\]
The probability of getting neither is
\[
1-\frac4{36}-\frac6{36}=\frac{26}{36}.
\]

For 5 to appear for the first time on roll $n$ before any 7,
the first $n-1$ rolls must be neither 5 nor 7, followed by a 5:
\[
P(\text{first qualifying result is 5 on roll $n$})
=
\left(\frac{26}{36}\right)^{n-1}\frac4{36}.
\]
Add over all possible values of $n$:
\[
\begin{aligned}
P(\text{5 before 7})
&=\sum_{n=1}^{\infty}
\left(\frac{26}{36}\right)^{n-1}\frac4{36}\\
&=\frac{4/36}{1-26/36}\\
&=\boxed{\frac25}.
\end{aligned}
\]

\subsection{Three distinguishable dice}

Roll a blue die $B$, a yellow die $Y$, and a red die $R$.
What is the probability that
\[
B<Y<R?
\]

For the values to be strictly increasing, they must first be
\emph{distinct}. Choose any three of the six faces. Once chosen,
exactly one arrangement places them in the order $B<Y<R$:
\[
\#\{\text{favorable outcomes}\}=\binom63=20.
\]
There are $6^3=216$ possible ordered outcomes, so
\[
\boxed{
P(B<Y<R)=\frac{\binom63}{6^3}
=\frac{20}{216}
=\frac5{54}.
}
\]

\subsection{Example: rain and lateness}

Suppose it rains tomorrow with probability $0.7$. Joe is late with
probability $0.3$ on rainy days and $0.1$ on nonrainy days.

Let $R$ mean ``rain'' and $L$ mean ``Joe is late.'' By the law of
total probability,
\[
\begin{aligned}
P(L)
&=P(R)P(L\mid R)+P(R^c)P(L\mid R^c)\\
&=(0.7)(0.3)+(0.3)(0.1)\\
&=\boxed{0.24}.
\end{aligned}
\]

Now suppose Joe was \emph{early}, meaning he was not late. What is
the probability that it rained? First,
\[
P(L^c)=1-0.24=0.76,
\qquad
P(L^c\mid R)=1-0.3=0.7.
\]
Bayes' formula gives
\[
\boxed{
P(R\mid L^c)
=\frac{P(R)P(L^c\mid R)}{P(L^c)}
=\frac{(0.7)(0.7)}{0.76}
=\frac{49}{76}.
}
\]

\subsection{Example: selecting a coin}

A box holds three coins: a two-headed coin, a fair coin, and a
biased coin that lands heads with probability $\tfrac34$.
One coin is chosen uniformly and flipped. Given that the result
is heads, what is the probability that the two-headed coin
was chosen?

Let $T$, $F$, and $D$ denote the three coin types, and let $H$
mean ``heads.'' Each coin is chosen with probability $\tfrac13$.
The law of total probability gives
\[
\begin{aligned}
P(H)
&=P(T)P(H\mid T)
 +P(F)P(H\mid F)
 +P(D)P(H\mid D)\\
&=\frac13(1)+\frac13\left(\frac12\right)
  +\frac13\left(\frac34\right)\\
&=\frac34.
\end{aligned}
\]
Therefore,
\[
\boxed{
P(T\mid H)
=\frac{P(T)P(H\mid T)}{P(H)}
=\frac{(1/3)(1)}{3/4}
=\frac49.
}
\]

\end{document}